% Solutions to the Chapter 4 exercises that are not code, from lectures/L04/appendix/b_exercises.md.
% Exercises 4.6 and 4.7 are Code and are not answered; see chapter.tex for why.

\section{Chapter 4: Kernel Modules}
\label{sol:sec:c4}
Answers to the exercises in \secref{c4:sec:exercises}. Exercises~\ref{c4:ex:6} and~\ref{c4:ex:7}
are Code, and are checked on the target as their Check yourself lines say.

% ------------------------------------------------------------------------------------------
\solution[fillaccent2]{c4:ex:1}{A module is not a program}{Recall}

\xpart{a} Nothing of the module's is running: its init function has returned and it has no thread
of its own. What is resident is its core memory: its code, including the exit function that
\code{rmmod} will call, its read-only data and parameter table, and its writable data, including
its \code{struct module}; with them, its directory under \file{/sys/module}. Its \code{__init} code
has already been freed. It waits to be called by something else.

\xpart{b} Because \code{insmod} is the process that called \code{init_module()}, and the kernel
runs your init function inside that call: \code{insmod} returns when init does. A \code{while (1)}
never returns, so neither does \code{insmod}. It sits in the system call, where a signal cannot
reach it, because a signal is acted on only on the way back to userspace; and the module stays in
the \code{Loading} state in \file{/proc/modules}, where \code{rmmod} refuses it. A CPU spins for
good, and the cure is a reboot. What you should have written is registration: a kernel thread
started with \code{kthread_run()} that loops until \code{kthread_should_stop()}, or a timer or a
work item if the work is periodic, started in init and stopped in exit.

\xpart{c} They are section attributes, and the linker and the loader act on them. \code{__init}
places a function in \code{.init.text}, and \code{__initdata} a variable in \code{.init.data}: for
code built into the kernel those sections are freed at the end of boot, and in a module the loader
puts them in separate init memory and frees it as soon as the init function has returned.
\code{__exit} places a function in \code{.exit.text}: a module keeps it, because \code{rmmod} will
call it, and a built-in one can never run it, so it is thrown away, at link time on most
architectures and with the init code on arm64 (\secref{c4:app:a1}). The consequence is that calling
an \code{__init} function after init, or keeping a pointer into \code{__initdata}, is a use after
free, and \code{modpost} warns of a section mismatch when ordinary code refers to init code.

\xpart{d} It is the kernel freeing what \code{__init} and \code{__initdata} set aside. Once the last
initcall has run, just before it starts \file{/init}, the kernel releases its whole init segment,
4,672 KB on this build, the stretch from \code{__init_begin} to \code{__init_end}
(\secref{c3:app:a8}). On arm64 most of that stretch is not \code{__init} code: 3.3 MB of it is the
relocation table the kernel applies to itself at boot, and the \code{__init} code and data are about
1 MB. The loader does the same for each module, on a small scale, when its init function returns
(Exercise~\ref{c4:ex:9}\textbf{f)}).

\xpart{e} \code{MODULE_LICENSE}, \code{MODULE_AUTHOR}, \code{MODULE_DESCRIPTION} and
\code{MODULE_VERSION}. \code{MODULE_LICENSE} is the one that is not optional: without it
\code{modpost} stops the build with \code{ERROR: modpost: missing MODULE_LICENSE() in}, followed by
the object's path, and no \code{.ko} is produced. A module that did load with no licence would be
treated as \code{unspecified}, which is not GPL-compatible: it would taint the kernel and see no
GPL-only symbol.

% ------------------------------------------------------------------------------------------
\solution[fillaccent2]{c4:ex:2}{Four commands}{Recall}

\xpart{a} \code{insmod} loads exactly the file it is given, with any parameters on its command
line, and resolves nothing. \code{modprobe} takes a module name, looks it up in the dependency file
under \code{/lib/modules/$(uname -r)}, and loads what it depends on first and then the module.
\code{depmod} scans the modules under \code{/lib/modules/$(uname -r)}, reads which symbols each one
needs and which it exports, and writes that dependency file. \code{rmmod} asks the kernel to unload
a module by name, which the kernel refuses while anything holds it.

\xpart{b} \code{insmod} resolves nothing and \code{modprobe} resolves dependencies, from a file
\code{depmod} wrote. The experiment: with \code{qa_export} not loaded, \code{insmod ./qa_user.ko}
fails; the kernel logs \code{qa_user: Unknown symbol qa_answer (err -2)} and \code{insmod} prints
\code{can't insert './qa_user.ko': unknown symbol in module or invalid parameter}
(\secref{c4:app:a4}). Then copy both modules under \code{/lib/modules/$(uname -r)/}, run
\code{depmod}, and run \code{modprobe qa_user}: \code{qa_export} is loaded first, then
\code{qa_user}, and \code{lsmod} shows both, with \code{qa_user} using \code{qa_export}.

\xpart{c} \file{modules.dep}, in \code{/lib/modules/$(uname -r)/}, which \code{depmod} generates;
\code{make modules_install} runs \code{depmod} for you (\secref{c3:app:b4}). BusyBox's simplified
\code{modprobe}, which is the one on this target, keeps its own version beside it,
\file{modules.dep.bb}, and generates that itself when it is missing.

\xpart{d} \file{/sys/module/qa_export/holders/}, which holds a link named \code{qa_user}.
\file{/sys/module/qa_export/refcnt} gives only the count, 1.

% ------------------------------------------------------------------------------------------
\solution[fillaccent3]{c4:ex:3}{Reading \code{vermagic}}{Hand calculation}

\xpart{a} \file{include/linux/vermagic.h} builds the string from six parts, each an assertion about
the kernel:
\begin{itemize}
  \item \code{6.12.30}: the kernel release it was built against.
  \item \code{SMP}: built with \code{CONFIG_SMP}, so its locking and per-CPU code assume more than
    one CPU.
  \item \code{preempt}: built for the preemptible-kernel model (\code{CONFIG_PREEMPT_BUILD}, which
    \code{CONFIG_PREEMPT} selects); a real-time kernel says \code{preempt_rt}, and a
    non-preemptible one says nothing.
  \item \code{mod_unload}: built with \code{CONFIG_MODULE_UNLOAD}, so modules can be unloaded and
    are reference-counted, which changes \code{struct module} and what the module's own code does.
  \item \code{modversions}: built with \code{CONFIG_MODVERSIONS}, so the module carries a checksum
    for every symbol it imports and the kernel checks them.
  \item \code{aarch64}: the architecture.
\end{itemize}

\xpart{b} Each of the four:
\begin{itemize}
  \item \textbf{Same source, no configuration change: it loads.} The \code{vermagic} string and
    every symbol checksum come out the same.
  \item \textbf{\code{\# CONFIG_SMP is not set}: it loads, because nothing changed.}
    \file{arch/arm64/Kconfig} declares \code{SMP} as \code{def_bool y}, with no prompt, so the line
    cannot take effect: \code{olddefconfig} keeps \code{CONFIG_SMP=y}, and the kernel is still SMP.
    Not even \file{ci/kernel.sh} says so, because it checks only the \code{CONFIG_X=} lines of a
    fragment, not the lines that turn something off.
  \item \textbf{6.12.30 to 6.12.31: it loads, unless one of the interfaces it uses changed.} With
    \code{modversions} the module has a table of checksums, and in that case the loader does not
    compare the release at all: \code{same_magic()} in \file{kernel/module/version.c} compares the
    two strings from the first space onwards. The rest of the string matches, so the decision falls
    to the checksums: of \code{module_layout}, which covers \code{struct module} and the structures
    around it, and of every symbol the module imports. A stable update rarely changes one; if it
    did, the loader prints \code{qa_hello: disagrees about version of symbol} and the symbol's name,
    and refuses. On a kernel without \code{CONFIG_MODVERSIONS} the release would be compared, and the
    module refused.
  \item \textbf{\code{PREEMPT_RT} instead of \code{PREEMPT}: it is refused.} The real-time kernel's
    string is \code{6.12.30 SMP preempt_rt mod_unload modversions aarch64}, and the comparison
    fails at \code{preempt} against \code{preempt_rt} (\secref{c4:app:a3}). It must: under
    \code{PREEMPT_RT} a \code{spinlock_t} is a sleeping lock with a different layout, and a module
    compiled for the other kernel would use every lock it touches wrongly.
\end{itemize}

\xpart{c} A checksum for every exported symbol, computed from the symbol's full type: its
prototype and, recursively, the layout of every structure it passes. The kernel records the
checksums of what it exports, and each module records, in its \code{__versions} section, the
checksum it was built against for each symbol it imports; the loader compares them one by one.
\code{vermagic} alone names a release and a few options, and cannot see a structure that changed
shape while the string stayed the same, through a configuration option that adds a field, a vendor
patch, or a rebuild with a different configuration and the same version number. The checksums check
the interface the module actually uses, which is also what lets the loader stop comparing the
release, so a module survives an update that did not touch anything it uses.

\xpart{d} Because the failure it prevents is silent memory corruption in the one address space
everything shares. A module built for other structure layouts reads and writes fields at the wrong
offsets, passes the wrong arguments and takes locks of the wrong size; nothing detects any of that
when it happens, and the damage surfaces later, somewhere else, as a crash or wrong data in code
that has nothing to do with the module (\secref{c3:app:a1}). A warning would scroll out of the log
and leave the corruption behind. A refusal costs a rebuild; the alternative costs a crash nobody can
trace, possibly in the field.

% ------------------------------------------------------------------------------------------
\solution[fillaccent3]{c4:ex:4}{The licence, at two stages}{Hand calculation}

\xpart{a} The build fails at \code{modpost}, the step that runs after the objects are linked and
before the \code{.ko} is finished. It sees \code{qa_export}'s GPL-only export and the proprietary
module's use of it, and stops with this, naming whatever your module is called:
\begin{console}
ERROR: modpost: GPL-incompatible module qa_proprietary.ko uses GPL-only symbol 'qa_answer'
\end{console}
\noindent No \code{.ko} is produced.

\xpart{b} The build still fails at \code{modpost}, but for an unresolved symbol:
\begin{console}
ERROR: modpost: "qa_answer" [.../qa_proprietary.ko] undefined!
\end{console}
\noindent The message is different because \code{modpost} can only apply the licence rule to an
export it has seen. Here \code{qa_answer} is in neither the kernel's \file{Module.symvers} nor any
module in the build, so as far as \code{modpost} knows the symbol does not exist, let alone that it
is GPL-only. Handing it the exporter's \file{Module.symvers} through \code{KBUILD_EXTRA_SYMBOLS}
brings back the message of \textbf{a)}.

\xpart{c} Whenever \code{modpost} could not see that the symbol is GPL-only. Either the module was
built as in \textbf{b)} with \code{KBUILD_MODPOST_WARN=1}, which turns the unresolved symbol into a
warning and lets the \code{.ko} be written; or it was built against a kernel, or an exporter's
\file{Module.symvers}, in which the symbol was plain \code{EXPORT_SYMBOL}, and is then loaded where
it has since become \code{EXPORT_SYMBOL_GPL}. The second is the one that happens on kernel upgrades.

\xpart{d} The kernel logs \code{qa_proprietary: Unknown symbol qa_answer (err -2)}, and BusyBox's
\code{insmod} prints \code{unknown symbol in module or invalid parameter}. No licence is mentioned
because the loader never refuses on grounds of licence. It checks the licence first, and marks the
module as proprietary; the symbol lookup then runs with \code{gplok} false, and
\code{find_exported_symbol_in_section()} skips every GPL-only table before searching it
(\secref{c4:app:a5}). The lookup fails exactly as it would for a symbol that is not exported at all,
with \code{-ENOENT}, and the loader reports, accurately, what it found: nothing.

\xpart{e} It arrives last, at load time on the target, possibly in the field after an upgrade,
rather than at build time on your own machine. It is word for word the message for a missing
dependency or a typo, so it sends you after the wrong causes. And the symbol is plainly in
\file{/proc/kallsyms}, so the message looks false. Diagnosing it requires already knowing the
mechanism. The first message is early and exact, the second early and at least says what is
missing; this one is late and points the wrong way.

% ------------------------------------------------------------------------------------------
\solution[fillaccent4]{c4:ex:5}{Choosing a log level}{Design}

\xpart{a} \code{pr_info}, or \code{dev_info} once there is a device: one line, once per device,
that says it is there and what it found is useful to whoever reads the log of a system that
misbehaves. The alternative, \code{pr_debug}, is what many maintainers prefer, on the grounds that a
working driver should be silent; that argument rests on a log that holds only what needs attention.

\xpart{b} \code{pr_warn}: something unexpected that the driver handled. The device misbehaved and
somebody may want to know, but the driver carries on. If it can happen on every sample, use
\code{pr_warn_ratelimited}, or it floods the log.

\xpart{c} \code{pr_err}: an error the driver could not handle. The probe fails and the device will
not work, so say why, with the error code.

\xpart{d} \code{pr_debug}: only useful during development, compiled to nothing unless dynamic debug
is built in, and silent until someone switches it on even then.

\xpart{e} \code{pr_warn}: the system works, but not as the device tree describes it, and the device
tree is what somebody has to fix.

\xpart{f} \textbf{The binary:} a \code{pr_info} you comment out is gone from the build, and getting
it back means editing and rebuilding. A \code{pr_debug} stays in the binary at the cost of a small
descriptor per call site and a branch that is patched to a no-op until enabled, and with
\code{CONFIG_DYNAMIC_DEBUG} off it compiles to nothing at all. \textbf{The field:} on a shipped
kernel with dynamic debug, those call sites can be switched on at run time on the machine that has
the problem, for one module, file, function or line, with \code{echo 'module qa_hello +p' >
/sys/kernel/debug/dynamic_debug/control}, without rebuilding or rebooting. You get the debugging
output you wrote, from the system that needs it.

% ------------------------------------------------------------------------------------------
\solution[fillaccent4]{c4:ex:8}{A test that cannot fail}{Design}

\xpart{a} 0. The kernel did refuse, \code{delete_module()} failing with \code{EWOULDBLOCK}, which is
\code{EAGAIN}, because \code{qa_export} has a user; and \code{rmmod} said so, \code{rmmod: remove
'qa_export': Resource temporarily unavailable}. But this target's BusyBox is built with its
simplified module tools, whose \code{rmmod} reports the failure and then exits with success
whatever happened.

\xpart{b} The obvious test is \code{if rmmod qa_export; then} fail. It sees 0 on both kernels, so
it concludes that the module was unloaded on both: a failure reported against the kernel that
correctly refused, and the same failure against the one that did not. Its verdict does not depend on
the kernel at all; what it measures is BusyBox. The same exit status, checked by a test that
expects an unload to succeed, passes on both, whether anything was unloaded or not, and that is
the test that cannot fail.

\xpart{c} Assert on the state you care about, not on a tool's report of it; and do not trust a
test you have never seen fail. The lab's check does both: it looks at \code{lsmod}, the state, and
the refusal it expects is one it can observe. From \secref{c2:app:a8}: BusyBox's \code{find} has no
\code{-printf}. A check written on a laptop with GNU \code{find}, which decides that a module is
gone because \code{find /sys/module -name qa_export -printf x} prints nothing, passes on the target
whether the module is there or not, because BusyBox's \code{find} rejects the option and prints
nothing on its standard output either way. It is the same class of problem: the check tests the
tool, which behaves differently on the machine it was written on and the one it runs on.

% ------------------------------------------------------------------------------------------
\solution[fillaccent]{c4:ex:9}{How big is a module}{Cross-check}

\textbf{a)}, \textbf{b)} and \textbf{d)} are measurements of your own module, and the course
publishes no figures for them; what follows is what decides each number.

\xpart{c} The prediction is the sum from \textbf{b)}: the module's code, its read-only data and
strings, its parameter table, the \code{.note} sections, and its writable data, including the
\code{.gnu.linkonce.this_module} section that holds its \code{struct module}, which on this kernel is
over a kilobyte by itself.

\xpart{e} Each of the six:
\begin{itemize}
  \item \textbf{The factor} is about five by the exercise's own figure, and the rest of this part is
    why: the loader rounds the module up in whole pages, not in bytes.
  \item \textbf{A multiple of 4,096}, the page size, because \code{module_memory_alloc()} rounds each
    memory type up to whole pages with \code{PAGE_ALIGN} before allocating it. The count is the
    number of core memory types the module has anything in, one page each as long as none of them
    holds more than 4,096 bytes, which for a module this small is at most four.
  \item \textbf{What would have to differ:} one module has something in a memory type the other has
    nothing in, or more than a page in one type. Between \code{qa_hello} and \code{qa_export} the
    usual difference is \code{qa_hello}'s \code{pr_debug}. With \code{CONFIG_JUMP_LABEL}, which this
    kernel has, each \code{pr_debug} call site has an entry in \code{__jump_table}, which the loader
    treats as read-only after init, a type of its own; a \code{qa_export} with no \code{pr_debug}
    has nothing of that type and needs a page fewer. Exporting does not add a type: the export
    tables are read-only data, which both modules already have.
  \item \textbf{The groups:} one per memory type the module uses. Sections of a type sit together at
    small offsets: code (\code{.text}, \code{.exit.text}), read-only data (\code{.rodata},
    \code{__param}, the \code{.note} sections), read-only after init (\code{__jump_table}), and
    writable data (\code{.data}, \code{.bss}, \code{.gnu.linkonce.this_module}). The init sections
    are listed too, at addresses that have already been freed. Between groups the addresses jump,
    because each type is a separate allocation from the module area, and vmalloc leaves an unmapped
    guard page after every allocation; where two sit next to each other, a one-page group is
    followed by the next two pages further on. A module is not one block but one allocation per
    memory type, each a whole number of pages.
  \item \textbf{\code{enum mod_mem_type}} has seven entries before \code{MOD_MEM_NUM_TYPES}: four
    core types, \code{MOD_TEXT}, \code{MOD_DATA}, \code{MOD_RODATA} and \code{MOD_RO_AFTER_INIT},
    and three init types. So the core size of a small module is 4,096 times the number of core types
    it uses, four pages at most. The allocation is \code{ptr = execmem_alloc(execmem_type, size);}
    in \code{module_memory_alloc()}, with \code{size} set by \code{PAGE_ALIGN} two lines earlier.
  \item \textbf{Why:} permissions are set per page. Code must be executable and read-only, read-only
    data neither writable nor executable, writable data not executable, and read-only-after-init
    data writable until init finishes and read-only after. A page holding two of them would need
    two permissions at once, which would leave code writable or data executable, exactly what
    \code{CONFIG_STRICT_MODULE_RWX} forbids; the loader sets them type by type in
    \file{kernel/module/strict_rwx.c}. The kernel's own image is laid out in separately aligned
    segments for the same reason (Exercise~\ref{c3:ex:9}).
\end{itemize}

\xpart{f} The init sections held \code{qa_hello_init} and anything marked \code{__initdata}. Once
the init function has returned, the loader frees them and sets the size of each init type to zero,
and \code{initsize} is the sum of those sizes. It is Exercise~\ref{c4:ex:1}\textbf{d)} in miniature:
the same annotation and the same freeing, done for one module by the loader after its init rather
than once by the kernel at the end of boot.

\xpart{g} The \code{coresize}, which is also the size in \file{/proc/modules}: ``the driver
occupies this many pages of kernel memory while it is loaded''. It is what the kernel actually
allocated and keeps; the file includes symbol tables and metadata that are never loaded, and the
sum of the sections leaves out the rounding to pages. With a 64 KB \code{kmalloc} in init, the
answer is the \code{coresize} plus 65,536 bytes: 64 KiB is above the largest \code{kmalloc} cache,
which is 8 KiB with 4 KiB pages, so it comes straight from the page allocator as sixteen pages, with
no rounding. None of the three numbers changes, apart from the few bytes of code the call adds,
because all three describe the module's own image: memory it allocates at run time belongs to the
allocator, and appears nowhere under \file{/sys/module}.
